\section*{第 \ref{ch:g6:decomposers-and-soil} 章——分解者与土壤}

\begin{solution}{exo:g6:decomposers-and-soil:1}
矿物颗粒（岩石碎屑）、各个分解阶段的有机遗体，
水分和空气，以及其中的生\omterm{def:g5:biodiversity:species}{物种}群。
深色的一层是\omterm{def:g6:decomposers-and-soil:soil}{腐殖质}。
\end{solution}

\begin{solution}{exo:g6:decomposers-and-soil:2}
以死亡\omterm{def:g6:plants-are-producers:matter}{有机物}为食——落叶、枯木、
粪便、尸体。可见的有：蚯蚓、潮虫、
马陆、跳虫（或粪甲虫、蝇幼虫）。
以菌丝毡为食者：\omterm{def:g6:decomposers-and-soil:decomposers}{真菌}。
\end{solution}

\begin{solution}{exo:g6:decomposers-and-soil:3}
还原为\omterm{def:g6:plants-are-producers:matter}{无机物}——先成为\omterm{def:g6:decomposers-and-soil:soil}{腐殖质}，再变为溶解的\omterm{prop:g4:what-plants-need:needs}{无机盐}，
并把气体释放到空气中。这恰与\omterm{def:g5:food-webs:roles}{生产者}的工作相反，
\omterm{def:g5:food-webs:roles}{生产者}把\omterm{def:g6:plants-are-producers:matter}{无机物}制成\omterm{def:g6:plants-are-producers:matter}{有机物}。
\end{solution}

\begin{solution}{exo:g6:decomposers-and-soil:4}
十一月：整片落下。冬季：被\omterm{def:g6:decomposers-and-soil:decomposers}{真菌}软化、分解，
被跳虫和潮虫啃蚀，只剩叶脉骨架。春季：
被蚯蚓拖入地下吞食，以碎屑形式排出。
夏季：变为黑色\omterm{def:g6:decomposers-and-soil:soil}{腐殖质}，其中的\omterm{def:g6:exploring-our-environment:components}{矿物质}溶于\omterm{def:g6:decomposers-and-soil:soil}{土壤}水，
被根吸收——或许正是这棵橡树自己的根。
\end{solution}

\begin{solution}{exo:g6:decomposers-and-soil:5}
因为\omterm{def:g6:decomposers-and-soil:decomposers}{分解者}吞食落叶的速度与每年落叶的速度一样快：
每一层落叶在下一层落满之前，
就已分解为\omterm{def:g6:decomposers-and-soil:soil}{腐殖质}和\omterm{def:g6:exploring-our-environment:components}{矿物质}。
\end{solution}

\begin{solution}{exo:g6:decomposers-and-soil:6}
温暖、潮湿、通气。冬季时分解工作几乎停止——
天气寒冷，这就是为什么落叶层会堆积得很厚，直到温暖
湿润的春季才将其清除。
\end{solution}

\begin{solution}{exo:g6:decomposers-and-soil:7}
\omterm{def:g5:food-webs:roles}{生产者} $\rightarrow$ \omterm{def:g5:food-webs:roles}{消费者}（\omterm{def:g1:plants-around-us:plant}{植物}被吃）；\omterm{def:g5:food-webs:roles}{生产者}
和\omterm{def:g5:food-webs:roles}{消费者} $\rightarrow$ 遗体、粪便、落叶
$\rightarrow$ \omterm{def:g6:decomposers-and-soil:decomposers}{分解者}（死亡\omterm{def:g6:exploring-our-environment:components}{生物}被吃掉）；
\omterm{def:g6:decomposers-and-soil:decomposers}{分解者} $\rightarrow$ \omterm{def:g6:decomposers-and-soil:soil}{土壤}水中的\omterm{prop:g4:what-plants-need:needs}{无机盐}
$\rightarrow$ 被\omterm{def:g5:food-webs:roles}{生产者}吸收：构成一个封闭的循环。
\end{solution}

\begin{solution}{exo:g6:decomposers-and-soil:8}
参与工作的有：蚯蚓、潮虫、跳虫、\omterm{def:g6:decomposers-and-soil:decomposers}{真菌}，以及看不见的
微\omterm{def:g6:exploring-our-environment:components}{生物}群。原料：厨房的有机废弃物。产物：
深色的\omterm{def:g6:decomposers-and-soil:soil}{腐殖质}，其中的\omterm{def:g6:exploring-our-environment:components}{矿物质}可供给花床。堆肥\omterm{def:g6:exploring-our-environment:conditions}{温度}较高，
因为这些工作者与所有\omterm{def:g5:food-webs:roles}{消费者}一样，会通过\omterm{def:g4:breathing:movements}{呼吸}作用消耗部分所吃的食物——
堆肥的热量正来自其“工作人员”自身的生命活动。
\end{solution}

\begin{solution}{exo:g6:decomposers-and-soil:9}
\omterm{def:g6:decomposers-and-soil:decomposers}{分解者}——把落叶拖入地下、吞食腐烂物质；
也是“耕耘者”：其洞穴使\omterm{def:g6:decomposers-and-soil:soil}{土壤}通气，
其排泄物使\omterm{def:g6:plants-are-producers:matter}{无机物}与\omterm{def:g6:plants-are-producers:matter}{有机物}充分混合。整块田地的
\omterm{def:g6:decomposers-and-soil:soil}{土壤}都会经过蚯蚓的肠道，这就是古代博物学家
誉之为“农田的造就者”的原因。
\end{solution}

\begin{solution}{exo:g6:decomposers-and-soil:10}
覆盖有落叶的花床：深色、疏松的\omterm{def:g6:decomposers-and-soil:soil}{腐殖质}；\omterm{def:g6:exploring-our-environment:components}{生物}普查数量丰富——蚯蚓、
潮虫、跳虫；地表凋落物逐季明显减少；
结论：物质循环运转良好，值得继续培育。
践踏过的小路：灰白、板结的尘土；几乎空无一物的托盘；
落下的东西都未腐烂；结论：贫瘠、缺氧的
\omterm{def:g6:decomposers-and-soil:soil}{土壤}，其中的工作者已经不见了。
\end{solution}

\begin{solution}{exo:g6:decomposers-and-soil:11}
缺氧。被密封且湿透的袋子隔绝了空气，
使需氧的工作者无法工作，只留下酸腐、无氧的腐烂过程。改进方法：
把材料从袋中取出，堆成疏松、湿润、经过翻动的堆——
保持\omterm{def:g6:exploring-our-environment:conditions}{湿度}，同时让空气进入。
\end{solution}

\begin{solution}{exo:g6:decomposers-and-soil:12}
堆积的是：死亡\omterm{def:g6:exploring-our-environment:components}{生物}——每年秋天的落叶层、每一代的遗体，
大量未被分解的\omterm{def:g6:plants-are-producers:matter}{有机物}。
耗尽的是：\omterm{def:g6:decomposers-and-soil:soil}{土壤}中的\omterm{prop:g4:what-plants-need:needs}{无机盐}，
一旦回归的箭头被切断，便不再得到补充；\omterm{def:g5:food-webs:roles}{生产者}吸收完
最后溶解的\omterm{def:g6:exploring-our-environment:components}{矿物质}后，便无法再合成\omterm{def:g6:plants-are-producers:matter}{有机物}，
无论光照多么充足。两个方面是同一条被切断的箭头——
分解作用——的两端。
\end{solution}

\begin{solution}{pb:g6:decomposers-and-soil:1}
\textbf{1.} 因此两个袋子只有一个不同之处——
网孔——任何质量差异都可归因于它：这是一个
公平实验（埋入土中）。

\textbf{2.} 看不见的\omterm{def:g6:decomposers-and-soil:soil}{土壤}\omterm{def:g6:exploring-our-environment:components}{生物}单独能否分解叶片——
粗网袋中既有可见的工作者，也有看不见的工作者；
细网袋则只保留看不见的\omterm{def:g6:exploring-our-environment:components}{生物}。

\textbf{3.} \omterm{def:g6:exploring-our-environment:conditions}{条件}是否决定分解工作：叶片相同、
网孔相同，但一处为干燥砾石，一处为潮湿河岸，这在野外检验
\cref{prop:g6:decomposers-and-soil:speed}中
关于潮湿\omterm{def:g6:exploring-our-environment:conditions}{条件}的论断。

\textbf{4.} 粗网袋损失最多，细网袋次之，密封袋几乎没有损失，
这一预测可直接依据各袋允许进入的
\omterm{def:g6:exploring-our-environment:components}{生物}和物质作出。

\textbf{5.} 粗网袋：全部工作者都参与了工作——分解最快。
细网袋：只有看不见的\omterm{def:g6:exploring-our-environment:components}{生物}参与工作——确有损失，但较慢。
密封袋：既无工作者，也无空气——几乎没有真正的分解，
只有酸腐、无氧的变质。

\textbf{6.} \omterm{def:g6:decomposers-and-soil:decomposers}{真菌}的菌丝侵入并软化叶片；
弹尾虫和鼠妇把叶片啃成网状；蚯蚓把碎片拖入地下，
吞食后排出加工过的碎屑；
看不见的工作者把每一片碎屑最终分解为\omterm{def:g6:decomposers-and-soil:soil}{腐殖质}和\omterm{prop:g4:what-plants-need:needs}{无机盐}。

\textbf{7.} 由此可见，它本身就是一支真正的\omterm{def:g6:decomposers-and-soil:decomposers}{分解者}队伍：
所有可见\omterm{def:g1:animals-around-us:animal}{动物}都被挡在外，袋内的重量仍在减少，
即分解仍在视觉以下的尺度上继续进行
。

\textbf{8.} 首要因素是\omterm{def:g6:exploring-our-environment:conditions}{湿度}——碎石排水、\omterm{def:g6:decomposers-and-soil:soil}{土壤}变干，
湿润堤岸所能容纳的工作种群随之减少；
干燥\omterm{def:g6:exploring-our-environment:conditions}{条件}使任何到来的工作队伍都无法开工。

\textbf{9.} 去向有二：进入\omterm{def:g6:decomposers-and-soil:soil}{土壤}与空气。一部分转化为\omterm{def:g6:decomposers-and-soil:soil}{腐殖质}，
并以溶解的\omterm{prop:g4:what-plants-need:needs}{无机盐}形式存在于堤岸\omterm{def:g6:decomposers-and-soil:soil}{土壤}水中；
另一部分被工作的\omterm{def:g6:decomposers-and-soil:decomposers}{分解者}通过\omterm{def:g4:breathing:movements}{呼吸}作用消耗，以气体释放，
损失的克数全部有去向，没有一克凭空消失。

\textbf{10.} 这并非巧合：埋入的叶片中的\omterm{prop:g4:what-plants-need:needs}{无机盐}
正是在那里释放出来，而荨麻是出了名的
喜肥\omterm{def:g1:plants-around-us:plant}{植物}。绿色条带是循环最后一步的可见标记：
分解作用为
\omterm{def:g5:food-webs:roles}{生产者}提供养分。

\textbf{11.} 垃圾填埋场的覆盖层相当于城市尺度上的密封袋：
无空气、无工作者，它剥夺了分解作用的
工作者和\omterm{def:g6:exploring-our-environment:conditions}{条件}，因此废物像袋中的污泥一样长期留存——
而同样的果皮若放在通风、潮湿、工作者众多的
堆肥中，数月便可回归\omterm{def:g6:decomposers-and-soil:soil}{土壤}。

\textbf{12.} 冬季的实验证明，落叶层是由\omterm{def:g6:decomposers-and-soil:soil}{土壤}中的\omterm{def:g6:decomposers-and-soil:decomposers}{分解者}清除的——
可见的工作者与看不见的工作者
共同参与，在温暖、潮湿、通风处最快。
它们的清除为\omterm{def:g6:decomposers-and-soil:soil}{土壤}的\omterm{prop:g4:what-plants-need:needs}{无机盐}储备提供补给，并经由这些\omterm{prop:g4:what-plants-need:needs}{无机盐}
惠及\omterm{def:g5:food-webs:roles}{生产者}：物质之轮
在树篱堤岸中转动。
\end{solution}
